\expandafter\gdef\csname SNside1\endcsname{\node[anchor=north west,inner sep=0pt,text width=48mm,font=\SideFont,align=left] (SNsidebox) at ([xshift=138mm,yshift=-402.91pt]current page.north west) {\textcolor{SNink}{\textbf{课程连接 · 套接字由操作系统管理}}\par\vspace{1.5mm}应用通过套接字接口使用传输服务，端口参与内核分用，进程号标识操作系统中的进程。一个服务器进程可接受多条连接，故不能把端口、进程和连接按一一对应关系死记。};\pgfextracty{\SNsideheight}{\pgfpointdiff{\pgfpointanchor{SNsidebox}{south}}{\pgfpointanchor{SNsidebox}{north}}}\typeout{SNHEIGHT:1:0:\the\SNsideheight}}
\expandafter\gdef\csname SNside2\endcsname{\node[anchor=north west,inner sep=0pt,text width=48mm,font=\SideFont,align=left] (SNsidebox) at ([xshift=138mm,yshift=-95.00pt]current page.north west) {\textcolor{SNink}{\textbf{不变量 · UDP长度包含8字节首部}}\par\vspace{1.5mm}应用数据12 B时UDP长度为20 B；普通IPv4还要把IP首部计入整个数据报上限。UDP自身保留报文边界，承载它的IP分片则必须先在目的IP层重组，不能把每个IP片都当一个UDP报文。};\pgfextracty{\SNsideheight}{\pgfpointdiff{\pgfpointanchor{SNsidebox}{south}}{\pgfpointanchor{SNsidebox}{north}}}\typeout{SNHEIGHT:2:2:\the\SNsideheight}}
\expandafter\gdef\csname SNside3\endcsname{\node[anchor=north west,inner sep=0pt,text width=48mm,font=\SideFont,align=left] (SNsidebox) at ([xshift=138mm,yshift=-216.43pt]current page.north west) {\textcolor{SNink}{\textbf{错项排除 · 伪首部不是真的多发一层}}\par\vspace{1.5mm}校验借用源IP、目的IP、协议和长度，是为检查交付对象；它只参加计算。奇数字节补0也只补计算过程。因此计算线上长度和吞吐量时，不能再给UDP额外增加伪首部和求和补位。};\pgfextracty{\SNsideheight}{\pgfpointdiff{\pgfpointanchor{SNsidebox}{south}}{\pgfpointanchor{SNsidebox}{north}}}\typeout{SNHEIGHT:3:3:\the\SNsideheight}}
\expandafter\gdef\csname SNside4\endcsname{\node[anchor=north west,inner sep=0pt,text width=48mm,font=\SideFont,align=left] (SNsidebox) at ([xshift=138mm,yshift=-192.52pt]current page.north west) {\textcolor{SNink}{\textbf{最小反例 · 双向seq各有自己的坐标}}\par\vspace{1.5mm}甲已发2001～4000，乙回复seq=4001、ack=3001：乙的seq属于反方向，只确认甲到3000的是ack。把乙的seq拿来确认甲的数据，会错误释放仍在途的3001～4000。};\pgfextracty{\SNsideheight}{\pgfpointdiff{\pgfpointanchor{SNsidebox}{south}}{\pgfpointanchor{SNsidebox}{north}}}\typeout{SNHEIGHT:4:5:\the\SNsideheight}}
\expandafter\gdef\csname SNside5\endcsname{\node[anchor=north west,inner sep=0pt,text width=48mm,font=\SideFont,align=left] (SNsidebox) at ([xshift=138mm,yshift=-95.00pt]current page.north west) {\textcolor{SNink}{\textbf{课程连接 · 状态是事件触发的记录}}\par\vspace{1.5mm}连接状态可像操作系统进程状态一样按事件画转换：发FIN、收到其ACK、收到对方FIN分别触发不同动作。仅看到报文名字不足以决定两端都已关闭，还要确定报文在何时被哪端收到。};\pgfextracty{\SNsideheight}{\pgfpointdiff{\pgfpointanchor{SNsidebox}{south}}{\pgfpointanchor{SNsidebox}{north}}}\typeout{SNHEIGHT:5:7:\the\SNsideheight}}
\expandafter\gdef\csname SNside6\endcsname{\node[anchor=north west,inner sep=0pt,text width=48mm,font=\SideFont,align=left] (SNsidebox) at ([xshift=138mm,yshift=-95.00pt]current page.north west) {\textcolor{SNink}{\textbf{错项排除 · 半关闭仍允许反向数据}}\par\vspace{1.5mm}被动端收到FIN后进入CLOSE-WAIT，表示对方不会再发送，本端应用仍可继续向对方发送。把TCP全双工看成一根双向同时关闭的管道，会错误断言收到FIN后所有方向的数据都不能发。};\pgfextracty{\SNsideheight}{\pgfpointdiff{\pgfpointanchor{SNsidebox}{south}}{\pgfpointanchor{SNsidebox}{north}}}\typeout{SNHEIGHT:6:8:\the\SNsideheight}}
\expandafter\gdef\csname SNside7\endcsname{\node[anchor=north west,inner sep=0pt,text width=48mm,font=\SideFont,align=left] (SNsidebox) at ([xshift=138mm,yshift=-95.00pt]current page.north west) {\textcolor{SNink}{\textbf{最小反例 · 首次加两次还差一次}}\par\vspace{1.5mm}先收到把确认号推进到3001的正常ACK，再收到两份ack=3001，重复计数只有2。第三份重复ACK到达才满足传统快重传阈值；同号ACK共有4份中的第一份不能混算成重复。\par\smallskip\begin{tikzpicture}[x=1mm,y=1mm,>=Stealth,node distance=5mm]
\node[draw=SNink!50,fill=SNblue,rounded corners=1mm,text width=34mm,align=center,inner sep=1.2mm,font=\fontsize{7.5}{9}\selectfont] (m0) at (0,0) {首ACK};
\node[draw=SNink!50,fill=SNblue,rounded corners=1mm,text width=34mm,align=center,inner sep=1.2mm,font=\fontsize{7.5}{9}\selectfont] (m1) at (0,-11) {重复1};
\node[draw=SNink!50,fill=SNblue,rounded corners=1mm,text width=34mm,align=center,inner sep=1.2mm,font=\fontsize{7.5}{9}\selectfont] (m2) at (0,-22) {重复2};
\node[draw=SNink!50,fill=SNblue,rounded corners=1mm,text width=34mm,align=center,inner sep=1.2mm,font=\fontsize{7.5}{9}\selectfont] (m3) at (0,-33) {尚未触发};
\draw[->,SNink] (m0) -- (m1);
\draw[->,SNink] (m1) -- (m2);
\draw[->,SNink] (m2) -- (m3);
\end{tikzpicture}
\par{\fontsize{7.5}{9}\selectfont 三个同号ACK仅含两份重复ACK}};\pgfextracty{\SNsideheight}{\pgfpointdiff{\pgfpointanchor{SNsidebox}{south}}{\pgfpointanchor{SNsidebox}{north}}}\typeout{SNHEIGHT:7:10:\the\SNsideheight}}
\expandafter\gdef\csname SNside8\endcsname{\node[anchor=north west,inner sep=0pt,text width=48mm,font=\SideFont,align=left] (SNsidebox) at ([xshift=138mm,yshift=-95.00pt]current page.north west) {\textcolor{SNink}{\textbf{不变量 · 按ACK事件依次更新四账}}\par\vspace{1.5mm}先更新累计确认，再核对已发未确认区间，随后按拥塞算法改变cwnd并读当前rwnd，最后取小扣在途。若只把两个窗口数字取小就回答还能发多少，会在所有有在途数据的题目中多算。};\pgfextracty{\SNsideheight}{\pgfpointdiff{\pgfpointanchor{SNsidebox}{south}}{\pgfpointanchor{SNsidebox}{north}}}\typeout{SNHEIGHT:8:16:\the\SNsideheight}}
\expandafter\gdef\csname SNside9\endcsname{\node[anchor=north west,inner sep=0pt,text width=48mm,font=\SideFont,align=left] (SNsidebox) at ([xshift=138mm,yshift=-240.29pt]current page.north west) {\textcolor{SNink}{\textbf{不变量 · 按ACK事件依次更新四账}}\par\vspace{1.5mm}先更新累计确认，再核对已发未确认区间，随后按拥塞算法改变cwnd并读当前rwnd，最后取小扣在途。若只把两个窗口数字取小就回答还能发多少，会在所有有在途数据的题目中多算。};\pgfextracty{\SNsideheight}{\pgfpointdiff{\pgfpointanchor{SNsidebox}{south}}{\pgfpointanchor{SNsidebox}{north}}}\typeout{SNHEIGHT:9:16:\the\SNsideheight}}
\expandafter\gdef\csname SNside10\endcsname{\node[anchor=north west,inner sep=0pt,text width=48mm,font=\SideFont,align=left] (SNsidebox) at ([xshift=138mm,yshift=-95.00pt]current page.north west) {\textcolor{SNink}{\textbf{不变量 · 按ACK事件依次更新四账}}\par\vspace{1.5mm}先更新累计确认，再核对已发未确认区间，随后按拥塞算法改变cwnd并读当前rwnd，最后取小扣在途。若只把两个窗口数字取小就回答还能发多少，会在所有有在途数据的题目中多算。};\pgfextracty{\SNsideheight}{\pgfpointdiff{\pgfpointanchor{SNsidebox}{south}}{\pgfpointanchor{SNsidebox}{north}}}\typeout{SNHEIGHT:10:16:\the\SNsideheight}}
\expandafter\gdef\csname SNside11\endcsname{\node[anchor=north west,inner sep=0pt,text width=48mm,font=\SideFont,align=left] (SNsidebox) at ([xshift=138mm,yshift=-95.00pt]current page.north west) {\textcolor{SNink}{\textbf{不变量 · 按ACK事件依次更新四账}}\par\vspace{1.5mm}先更新累计确认，再核对已发未确认区间，随后按拥塞算法改变cwnd并读当前rwnd，最后取小扣在途。若只把两个窗口数字取小就回答还能发多少，会在所有有在途数据的题目中多算。};\pgfextracty{\SNsideheight}{\pgfpointdiff{\pgfpointanchor{SNsidebox}{south}}{\pgfpointanchor{SNsidebox}{north}}}\typeout{SNHEIGHT:11:16:\the\SNsideheight}}
\expandafter\gdef\csname SNside12\endcsname{\node[anchor=north west,inner sep=0pt,text width=48mm,font=\SideFont,align=left] (SNsidebox) at ([xshift=138mm,yshift=-95.00pt]current page.north west) {\textcolor{SNink}{\textbf{不变量 · 按ACK事件依次更新四账}}\par\vspace{1.5mm}先更新累计确认，再核对已发未确认区间，随后按拥塞算法改变cwnd并读当前rwnd，最后取小扣在途。若只把两个窗口数字取小就回答还能发多少，会在所有有在途数据的题目中多算。};\pgfextracty{\SNsideheight}{\pgfpointdiff{\pgfpointanchor{SNsidebox}{south}}{\pgfpointanchor{SNsidebox}{north}}}\typeout{SNHEIGHT:12:16:\the\SNsideheight}}
\expandafter\gdef\csname SNside13\endcsname{\node[anchor=north west,inner sep=0pt,text width=48mm,font=\SideFont,align=left] (SNsidebox) at ([xshift=138mm,yshift=-95.00pt]current page.north west) {\textcolor{SNink}{\textbf{不变量 · 按ACK事件依次更新四账}}\par\vspace{1.5mm}先更新累计确认，再核对已发未确认区间，随后按拥塞算法改变cwnd并读当前rwnd，最后取小扣在途。若只把两个窗口数字取小就回答还能发多少，会在所有有在途数据的题目中多算。};\pgfextracty{\SNsideheight}{\pgfpointdiff{\pgfpointanchor{SNsidebox}{south}}{\pgfpointanchor{SNsidebox}{north}}}\typeout{SNHEIGHT:13:16:\the\SNsideheight}}
\expandafter\gdef\csname SNside14\endcsname{\node[anchor=north west,inner sep=0pt,text width=48mm,font=\SideFont,align=left] (SNsidebox) at ([xshift=138mm,yshift=-95.00pt]current page.north west) {\textcolor{SNink}{\textbf{课程连接 · 两层缓冲分开看}}\par\vspace{1.5mm}应用缓冲与内核TCP收发缓冲不同；对端应用读取速度会影响接收窗口。};\pgfextracty{\SNsideheight}{\pgfpointdiff{\pgfpointanchor{SNsidebox}{south}}{\pgfpointanchor{SNsidebox}{north}}}\typeout{SNHEIGHT:14:17:\the\SNsideheight}}
\expandafter\gdef\csname SNside15\endcsname{\node[anchor=north west,inner sep=0pt,text width=48mm,font=\SideFont,align=left] (SNsidebox) at ([xshift=138mm,yshift=-95.00pt]current page.north west) {\textcolor{SNink}{\textbf{计数 · 原ACK加三个重复}}\par\vspace{1.5mm}首次ACK623推进确认；再收到三个相同的重复ACK才达经典快重传阈值。};\pgfextracty{\SNsideheight}{\pgfpointdiff{\pgfpointanchor{SNsidebox}{south}}{\pgfpointanchor{SNsidebox}{north}}}\typeout{SNHEIGHT:15:20:\the\SNsideheight}}
\expandafter\gdef\csname SNside16\endcsname{\node[anchor=north west,inner sep=0pt,text width=48mm,font=\SideFont,align=left] (SNsidebox) at ([xshift=138mm,yshift=-95.00pt]current page.north west) {\textcolor{SNink}{\textbf{计数 · 原ACK加三个重复}}\par\vspace{1.5mm}首次ACK623推进确认；再收到三个相同的重复ACK才达经典快重传阈值。};\pgfextracty{\SNsideheight}{\pgfpointdiff{\pgfpointanchor{SNsidebox}{south}}{\pgfpointanchor{SNsidebox}{north}}}\typeout{SNHEIGHT:16:20:\the\SNsideheight}}
\AddToShipoutPictureFG{\begin{tikzpicture}[remember picture,overlay]\ifcsname SNside\thepage\endcsname\csname SNside\thepage\endcsname\fi\end{tikzpicture}}